% Submission-style pre-study manuscript, 12 September 2026.
% Core claim/history register: research_record.md. Preserve claims and evidence.
% Earlier manuscript retained verbatim in proposal_support/archive/before_2027_submit_writeup/.
\documentclass[11pt,a4paper]{article}
\usepackage[margin=1in]{geometry}
\usepackage{fontspec}
\setmainfont{Times New Roman}
\usepackage{amsmath,amssymb,amsthm,booktabs,enumitem}
\usepackage[hidelinks]{hyperref}
\setlength{\emergencystretch}{2em}
\setlist{itemsep=2pt,topsep=3pt}
\newtheorem{theorem}{Theorem}[section]
\newtheorem{proposition}[theorem]{Proposition}
\newtheorem{corollary}[theorem]{Corollary}
\theoremstyle{definition}
\newtheorem{definition}[theorem]{Definition}
\theoremstyle{remark}
\newtheorem{remark}[theorem]{Remark}
\newcommand{\R}{\mathbb R}
\newcommand{\E}{\mathcal E}
\newcommand{\norm}[1]{\left\|#1\right\|_\infty}
\makeatletter
\let\tableinput\@@input
\makeatother
\input{proposal_support/lagrange_evidence/pilot_values.tex}
\title{Lie-Symmetry Reduction and Taylor--PINN Approximation\\
of a Lagrangian Shallow-Water Equation}
\author{Carl\\WKU MATH 4890}
\date{Pre-study manuscript for an eight-week undergraduate project in 2027\\
Preparatory evidence: 12 September 2026}
\hypersetup{pdftitle={Lie-Symmetry Reduction and Taylor-PINN Approximation of a Lagrangian Shallow-Water Equation},pdfauthor={Carl}}
\begin{document}
\maketitle

\begin{abstract}
We study a selected invariant family of the third-order Lagrangian master equation introduced by Paliathanasis, Leon and Leach (2026). For cubic sediment-flux exponent, the invariant $\sigma=st^2$ reduces the printed equation to a regular third-order profile ODE away from an explicit leading-coefficient zero set. We derive the reduction directly, prove an exact residual-lifting identity, and give a conditional a-posteriori continuation and error estimate. Exact rational interval arithmetic establishes a state Lipschitz constant of 15 on a specified slope--curvature box. These results support a comparison of Taylor polynomials, adaptive integration and hard-initial-data Taylor--PINN profiles. Preparatory symbolic checks, a classical numerical pilot and two short CPU neural runs are reported separately from the planned 2027 multi-seed study. The neural runs show different rankings for profile error and residual, so neither higher Taylor order nor lower sampled residual establishes superior approximation. A chain-rule audit also distinguishes the printed master equation from the usual Eulerian sediment-transport model: an exact solution of the former need not satisfy the latter. The contribution is a source-specific, reproducible reduction and verification study, not a new symmetry classification, general PDE solver or established neural performance advantage.
\end{abstract}

\section{Introduction and Mathematical Problem}\label{sec:problem}
Lie symmetry reduction restricts a PDE to solutions unchanged by a selected continuous transformation. Its benefit is exact removal of an independent variable; its limitation is restriction to a special solution family \cite{olver}. We use this classical operation before numerical approximation. The project is deliberately limited to one equation, one principal reduction and one small family of profile data.

Our primary source is Paliathanasis, Leon and Leach \cite{lagrange}. Its printed Equation (7), with dimensionless fixed parameters $g,\beta>0$, is
\begin{equation}
\E_m[\phi]:=
\partial_t\left(\phi_s\phi_{tt}-g\frac{\phi_{ss}}{\phi_s^2}\right)
-g\beta\partial_s\left(\phi_t^{m-1}\frac{\phi_{ts}}{\phi_s}\right)=0.
\label{eq:master}
\end{equation}
Here $\phi(t,s)$ is particle position, $t$ is time and $s$ is a fluid-mass label. On a regular positive branch,
\begin{equation}
x=\phi(t,s),\qquad u=\phi_t>0,\qquad h=\phi_s^{-1}>0.
\label{eq:map}
\end{equation}
The source parameter is $\beta=mA_g\xi$, not $A_g\xi$, for sediment flux $Q=A_gu^m$. We fix $m=3$, $g=1$ and $\beta=1/10$ in all principal numerical experiments.

\paragraph{Specific problem.}
Using the source's cubic scaling $X_5=t\partial_t-2s\partial_s$, seek
\begin{equation}
\phi(t,s)=F(\sigma),\qquad \sigma=st^2,\qquad
\sigma\in I=[1,3/2].
\label{eq:ansatz}
\end{equation}
The profile satisfies the reduced initial-value problem
\begin{equation}
a(\sigma,F')F'''+b(\sigma,F',F'')=0,\qquad
\bigl(F,F',F''\bigr)(1)=
\left(K,\frac K3,-\frac{2K}{9}(1+\delta)\right),
\quad K=\sqrt[3]{45},
\label{eq:ivp}
\end{equation}
with $\delta=0$ as an exact control and $\delta=0.01$ as a perturbed-profile test. Coefficients $a,b$ are derived in Section~\ref{sec:reduction}. The symbol $\delta$ changes curvature data, not the exponent $m$. These are profile initial data on the similarity coordinate, not arbitrary time-initial or boundary data for the full PDE.

For reconstruction define
\begin{equation}
\Omega=\{(t,s):1\le t\le5/4,\ 1\le st^2\le3/2\}.
\label{eq:domain}
\end{equation}
For an approximation $\widehat F$, the quantities to distinguish are
\begin{equation}
e_F=\|\widehat F-F_{\rm ref}\|_{\infty,I},\quad
e_Y=\|(\widehat F,\widehat F',\widehat F'')-Y_{\rm ref}\|_{\infty,I},
\quad
\rho=\|\E_3[\widehat F(st^2)]\|_{\infty,\Omega}.
\label{eq:metrics}
\end{equation}
A grid maximum is denoted explicitly as sampled; it is not substituted for any supremum in a theorem. The research question is whether a small Taylor--PINN trial improves a declared error--cost trade-off over finite Taylor polynomials and classical integration, while maintaining regularity. The answer may be negative.

\paragraph{Evidence and scope.}
The selected scaling and power-law family are prior work. We independently verify the formulas needed here, rather than adopting every printed reduction or the completeness of the source's classification. The main target remains the printed equation~\eqref{eq:master}, as approved for this project. Physical-model correction, shock treatment, maximal symmetry discovery and general PDE well-posedness are outside the eight-week scope.

\section{Relationship to Prior Work}\label{sec:prior}
Nagesh et al.~\cite{taylor} construct Taylor-model PINNs by combining Taylor expansions with neural remainders and deriving remainder dynamics. Their Section 3.1 and Theorems 4--6 provide direct methodological precedent. Their implementation works with an ODE state and avoids the higher neural derivatives used in our scalar third-order residual. Consequently our construction is not a reproduction of their complete method, and their error guarantees do not automatically transfer. Their reported longer-horizon deterioration and additional per-epoch cost also preclude an assumed advantage.

Residual collocation follows the standard PINN framework \cite{pinn}. Akhound-Sadegh et al.~\cite{liepinn} introduce symmetry-informed losses and discuss useful, zero and redundant generator contributions in specific experiments. Shah et al.~\cite{shah} instead transform collocation points for Burgers-type problems. Neither source supports the assertion that all symmetry-informed learning fails. Our distinction is narrower: an invariant representation is fixed analytically before a profile is learned.

McGreivy and Hakim \cite{baselines} motivate efficient classical baselines, matched-accuracy comparisons and transparent training costs. They do not require exclusively closed-form references or prohibit numerical reference solutions. We therefore separate an exact control, a highly accurate cross-checked reference, and a practically tuned integration baseline. Combining established ingredients is not, by itself, evidence of novelty.

\section{Model Audit and Bed Reconstruction}\label{sec:audit}
For comparison only, the source's Eulerian equations on $h>0$ can be written
\begin{equation}
h_t+(hu)_x=0,\qquad
u_t+uu_x+g(h+B)_x=0,\qquad
B_t+\beta u^{m-1}u_x=0.
\label{eq:euler}
\end{equation}
The primitive momentum equation is equivalent to the printed conservative momentum balance when mass conservation and $h>0$ hold. Let the usual bed pullback be $\widetilde B(t,s)=B(t,\phi(t,s))$. The chain rule gives
\begin{equation}
\partial_x=\phi_s^{-1}\partial_s,\qquad
(\partial_t)_x=(\partial_t)_s-\frac{\phi_t}{\phi_s}\partial_s.
\label{eq:chain}
\end{equation}
Thus the physical Exner equation becomes
\begin{equation}
\widetilde B_t-\frac{\phi_t}{\phi_s}\widetilde B_s+
\beta\phi_t^{m-1}\frac{\phi_{ts}}{\phi_s}=0.
\label{eq:physicalbed}
\end{equation}
The transport term is absent from printed Equation (6) of \cite{lagrange}. This is a precise inconsistency under the stated pullback, not a conclusion about an unspecified alternative transformed bed variable.

Set
\begin{equation}
A=\phi_s\phi_{tt}-g\phi_{ss}/\phi_s^2,\qquad
q=\phi_t^{m-1}\phi_{ts}/\phi_s.
\label{eq:Aq}
\end{equation}
The printed momentum and bed equations specify
\begin{equation}
\widetilde B_s=-A/g,\qquad \widetilde B_t=-\beta q.
\label{eq:printedpair}
\end{equation}

\begin{proposition}[Reconstruction for the printed pair]\label{prop:bed}
Let $U$ be a simply connected open region and $\phi\in C^3(U)$ with $\phi_s>0$ and $\phi_t>0$. Then $\E_m[\phi]=0$ if and only if the one-form
\[
\omega=-\beta q\,dt-(A/g)\,ds
\]
is closed. In that case there exists a $C^2$ function $\widetilde B$ satisfying~\eqref{eq:printedpair}, unique up to an additive constant on $U$.
\end{proposition}
\begin{proof}
The coefficient of $dt\wedge ds$ in $d\omega$ is
$-A_t/g+\beta q_s=-\E_m/g$. Closedness is therefore exactly the master equation. A closed continuously differentiable one-form on a simply connected region is exact, giving $d\widetilde B=\omega$. Two such potentials have zero differential and differ by a constant.
\end{proof}

\begin{remark}
This proposition reconstructs the printed pair, not~\eqref{eq:physicalbed}. Restoring transport yields the different compatibility equation
\[
\E_m[\phi]-\partial_s(\phi_tA/\phi_s)=0.
\]
Its use would require a separate model/reduction study. A small nonzero master residual also does not itself guarantee a path-independent approximate bed.
\end{remark}

\subsection{Exact Negative Control}\label{sec:negative}
For $m=3$, $g=\beta=1$, take
\begin{equation}
\phi=-\log(2-s-t),\qquad
\widetilde B=s-\frac{1}{2(2-s-t)^2},\qquad 0\le s+t\le1.
\label{eq:wave}
\end{equation}
Direct differentiation shows $\phi_s=\phi_t=(2-s-t)^{-1}>0$ and all three printed residuals, including~\eqref{eq:master}, vanish. In Eulerian coordinates this gives
\begin{equation}
h=e^{-x},\qquad u=e^x,\qquad
B=2-t-e^{-x}-\tfrac12e^{2x}.
\end{equation}
Mass and momentum residuals in~\eqref{eq:euler} are zero, but
\begin{equation}
B_t+u^2u_x=e^{3x}-1\not\equiv0.
\label{eq:counterexample}
\end{equation}
The verification script checks each expression symbolically. This negative result is retained throughout the study.

\section{Cubic Symmetry Reduction}\label{sec:reduction}
\begin{proposition}[Selected cubic symmetry]\label{prop:symmetry}
The transformation generated by $X_5=t\partial_t-2s\partial_s$ preserves~\eqref{eq:master} at $m=3$, wherever both transformed and original expressions are defined. Its invariant profiles on $t,s>0$ have the form~\eqref{eq:ansatz}.
\end{proposition}
\begin{proof}
The graph action is $(T_\epsilon\phi)(t,s)=\phi(e^{-\epsilon}t,e^{2\epsilon}s)$. The first term in~\eqref{eq:master} gains factor $e^{-\epsilon}$, and the second gains $e^{(2-m)\epsilon}$. At $m=3$ both factors agree, so
$\E_3[T_\epsilon\phi]=e^{-\epsilon}T_\epsilon(\E_3[\phi])$.
The invariant-surface condition $t\phi_t-2s\phi_s=0$ has characteristic invariant $st^2$, giving $\phi=F(st^2)$.
\end{proof}
We verify this particular source symmetry; no maximality statement is needed. The finite experimental domain~\eqref{eq:domain} is a restriction of the invariant family, not a claim that every group element preserves its boundary.

\begin{theorem}[Exact profile reduction and residual lift]\label{thm:reduction}
Let $F\in C^3(I)$, $r=F'>0$, $v=F''$, $t>0$ and $\sigma=st^2\in I$. Define
\begin{align}
a(\sigma,r)&=8(1-g\beta)\sigma^3r-\frac{2g\sigma}{r^2},\label{eq:a}\\
b(\sigma,r,v)&=4\sigma(1-4g\beta)r^2
+8\sigma^2(3-5g\beta)rv
+\left(8\sigma^3(1-g\beta)+\frac{4g\sigma}{r^3}\right)v^2.
\label{eq:b}
\end{align}
Then
\begin{equation}
\E_3[F(st^2)]=t^{-1}R[F],\qquad R[F]=a(\sigma,F')F'''+b(\sigma,F',F'').
\label{eq:lift}
\end{equation}
If $a\ne0$, the reduced equation is the regular state system
\begin{equation}
Y'=G(\sigma,Y),\qquad Y=(F,r,v),\qquad
G=(r,v,f),\qquad f=-b/a.
\label{eq:state}
\end{equation}
\end{theorem}
\begin{proof}
Differentiation gives
\[
\phi_s=t^2r,\quad \phi_t=2\sigma r/t,\quad
\phi_{ss}=t^4v,\quad \phi_{tt}=(2\sigma r+4\sigma^2v)/t^2,\quad
\phi_{ts}=2t(r+\sigma v).
\]
Hence $A=2\sigma r^2+4\sigma^2rv-gv/r^2$ and
$q=t^{-3}C$, where $C=8\sigma^2r(r+\sigma v)$.
At fixed $s$, $\partial_tA=2\sigma A'/t$; at fixed $t$,
$\partial_sq=C'/t$. Therefore
$t\E_3=2\sigma A'-g\beta C'$. Expanding and collecting the coefficient of
$v'=F'''$ yields exactly~\eqref{eq:a}--\eqref{eq:b}. Dividing by $a$ is legitimate only on the declared regular branch.
\end{proof}

The excluded set is explicit:
\begin{equation}
a=0\quad\Longleftrightarrow\quad
4(1-g\beta)\sigma^2r^3=g.
\label{eq:rankdrop}
\end{equation}
It must not be removed by cancellation or crossed without a new analysis.

\begin{proposition}[Positive power-law control]\label{prop:control}
If $g\beta<1/4$, then, for any constant $C_0$,
\[
F_0=C_0+K\sigma^{1/3},\qquad K^3=\frac{27g}{1-4g\beta},\quad K>0,
\]
satisfies $R[F_0]=0$, $F_0'>0$, and
\begin{equation}
a(\sigma,F_0')=2K\sigma^{7/3}>0\qquad(\sigma>0).
\label{eq:exacta}
\end{equation}
\end{proposition}
\begin{proof}
Substitution gives
$R[F_0]=-4(4K^3g\beta-K^3+27g)/(81K\sigma^{1/3})$,
which vanishes by the root relation. The slope is
$K\sigma^{-2/3}/3$. Substituting this slope into~\eqref{eq:a}
and using the same root relation gives~\eqref{eq:exacta}.
\end{proof}
The family and root relation occur in \cite[Section 3.4]{lagrange}.
Equation~\eqref{eq:exacta} records the regularity required for our chosen control. This smooth positive-domain calculation is not a shock construction.

\section{Conditional Error Control}\label{sec:theory}
All state norms in this section are infinity norms. The mathematical result is an application of classical ODE theory, not a new general error theorem.

\begin{proposition}[Conditional residual estimate]\label{prop:gronwall}
Let an exact solution $Y$ of~\eqref{eq:state} and a trial
$\widehat Y=(\widehat F,\widehat F',\widehat F'')$, $\widehat F\in C^3(I)$,
remain in a common region on which $G$ is uniformly $L$-Lipschitz in the state. Suppose
\[
\|\widehat Y(1)-Y(1)\|_\infty\le e_0,\qquad
|\widehat F'''-f(\sigma,\widehat F',\widehat F'')|\le\varepsilon
\quad\hbox{for every }\sigma\in I.
\]
Then, for $z=\sigma-1$,
\begin{equation}
\|\widehat Y(\sigma)-Y(\sigma)\|_\infty
\le e_0e^{Lz}+\varepsilon\frac{e^{Lz}-1}{L},
\label{eq:gronwall}
\end{equation}
with the quotient interpreted as $z$ when $L=0$.
\end{proposition}
\begin{proof}
The trial state satisfies $\widehat Y'=G(\sigma,\widehat Y)+(0,0,d)$ with $|d|\le\varepsilon$. Subtracting the integral equations yields
$\|\widehat Y-Y\|\le e_0+\int_1^\sigma(L\|\widehat Y-Y\|+\varepsilon)\,d\rho$.
The integral Gronwall inequality gives~\eqref{eq:gronwall}.
\end{proof}

\subsection{A Computable Regularity Box}
For the numerical parameters, let
\begin{equation}
\mathcal B=[17/20,5/4]\times[-17/20,-7/20]
\quad\hbox{for }(r,v),\qquad I=[1,3/2].
\label{eq:box}
\end{equation}
Here
\begin{equation}
f(\sigma,r,v)=
-\frac{2(3r^5+25\sigma r^4v+9\sigma^2r^3v^2+5v^2)}
{r(18\sigma^2r^3-5)}.
\label{eq:rationalf}
\end{equation}

\begin{proposition}[Verified box constants]\label{prop:box}
On $I\times\mathcal B$,
\[
\frac{24217}{7225}\le a\le\frac{5691}{200},\qquad
\|D_YG\|_\infty\le15.
\]
\end{proposition}
\begin{proof}
The derivative $a_r=8(1-g\beta)\sigma^3+4g\sigma/r^3$ is positive, and
$a_\sigma=24(1-g\beta)\sigma^2r-2g/r^2$ is positive throughout this box.
Evaluation at its two diagonal corners gives the stated endpoints.
For the Jacobian, its row sums are $1,1,|f_r|+|f_v|$.
Appendix~\ref{app:interval} gives the differentiated rational expressions and exact interval procedure. A Cartesian partition into $12^3$ rational boxes yields
\[
|f_r|+|f_v|
\le\frac{7150105171934045737}{506087956699582464}
<15.
\]
All interval operations use exact fractions, with strictly positive denominators. This is an upper enclosure, not a claim that the derivative maximum is attained at that value.
\end{proof}

\begin{theorem}[Conditional continuation certificate]\label{thm:certificate}
Let $\widehat F\in C^3(I)$, and suppose its slope--curvature pair lies inside $\mathcal B$, at infinity-distance at least $\eta>0$ from $\partial\mathcal B$ throughout $I$. Suppose its normalized defect is uniformly bounded by $\varepsilon$, and let the desired initial state satisfy
$\|Y_0-\widehat Y(1)\|_\infty\le e_0$. If
\begin{equation}
\Delta=e_0e^{15/2}+\frac{\varepsilon}{15}(e^{15/2}-1)<\eta,
\label{eq:certificate}
\end{equation}
then the IVP with initial data $Y_0$ has a unique solution on $I$,
its slope--curvature pair remains inside $\mathcal B$, and
$\|Y-\widehat Y\|_{\infty,I}\le\Delta$.
\end{theorem}
\begin{proof}
The strict inequality puts $Y_0$ in the regular region. Local existence and uniqueness follow from smoothness of~\eqref{eq:rationalf}. Until a hypothetical first exit of $(r,v)$ from $\mathcal B$, convexity of the box and Proposition~\ref{prop:box} permit Proposition~\ref{prop:gronwall} with $L=15$.
The resulting error is strictly smaller than the trial's boundary margin, contradicting an exit. A solution cannot terminate earlier inside this region: $r,v$ remain bounded, and $F'=r$ bounds $F$ on the finite interval. The usual ODE continuation argument therefore extends the solution to the right endpoint.
\end{proof}

The matched-initial-data factor is
$(e^{7.5}-1)/15\approx120.4695$. Thus even this verified constant can yield an unhelpful estimate. Neither a sampled residual nor a sampled margin discharges the theorem's uniform hypotheses. A failure of~\eqref{eq:certificate} is reported as inconclusive, not as proof that the approximation is wrong.

\begin{corollary}[Residual and field reconstruction]\label{cor:lift}
For a $C^3$ trial with $\widehat F'>0$ on $I$ and
$\Omega_*=\{(t,s):0<t_{\min}\le t\le t_{\max},\ st^2\in I\}$,
\[
\|\E_3[\widehat F(st^2)]\|_{\infty,\Omega_*}
\le t_{\min}^{-1}\|R[\widehat F]\|_{\infty,I}.
\]
If exact and trial slopes share a positive lower bound $r_{\min}$, then at the same labels $(t,s)$,
\[
|\widehat h-h|\le
\frac{|\widehat F'-F'|}{t_{\min}^2r_{\min}^2},\qquad
|\widehat u-u|\le\frac{2\max I}{t_{\min}}|\widehat F'-F'|.
\]
\end{corollary}
\begin{proof}
The first estimate is Theorem~\ref{thm:reduction}. The other two follow from
$h=(t^2F')^{-1}$ and $u=2\sigma F'/t$.
\end{proof}
These are neither fixed-Eulerian-location field errors nor bed errors, and they do not certify~\eqref{eq:euler}.

\section{Semi-Analytical Neural Method}\label{sec:method}
Let $z=\sigma-1$ and write the initial state as $(F_*,r_*,v_*)$.
The analytically determined Taylor coefficients are
\begin{align}
P_2(z)&=F_*+r_*z+\tfrac12v_*z^2,\\
P_4(z)&=P_2(z)+\tfrac16 f_*z^3+
\tfrac1{24}(f_\sigma+f_rv+f_vf)_*z^4.
\label{eq:taylor}
\end{align}
Smoothness of $f$ on the regular branch justifies these local derivatives. A finite polynomial comparison is not a claim of convergence of an infinite Taylor series on $I$.

Define the scalar trial
\begin{equation}
\widehat F_{\theta,N}=P_N(z)+z^{N+1}n_\theta(4z-1),
\qquad N\in\{2,4\},
\label{eq:trial}
\end{equation}
where $n_\theta$ has two tanh hidden layers of width 16 and a scalar linear output. It has 321 trainable parameters. Both trials enforce the three initial data exactly; $N=4$ also fixes the next two derivatives to their ODE values. The factors $z^3$ and $z^5$ change optimization conditioning as well as Taylor information. The comparison therefore concerns combined parameterizations, not the isolated causal effect of Taylor order.

For training points $\sigma_j$, the objective is
\begin{equation}
\mathcal L_N(\theta)=\frac1{M}\sum_{j=1}^M
\left[
\widehat F_{\theta,N}'''(\sigma_j)
+\frac{b(\sigma_j,\widehat F_{\theta,N}',\widehat F_{\theta,N}'')}
{a(\sigma_j,\widehat F_{\theta,N}')}
\right]^2,\qquad M=128.
\label{eq:loss}
\end{equation}
No reference solution labels enter this loss. All derivatives are taken from the single output $\widehat F$; independent auxiliary outputs cannot silently replace its derivatives. Both the normalized defect and raw residual $R$ are evaluated, since coefficient normalization changes their relative scales.

\paragraph{Regularity and initialization.}
The final neural layer is initialized to zero, so the initial trial is precisely $P_N$; hidden layers retain seeded random initialization. Computation uses float64 on one CPU thread. A sampled $F'\le0$, $|a|<10^{-6}$, nonfinite value or nonfinite gradient marks the run as failed. No denominator clipping is used. Such checks protect evaluated points, not the entire continuous interval.

\paragraph{Comparison protocol.}
The planned study compares $P_2$, $P_4$, DOP853 and the two neural trials for $\delta=0,0.01$. Neural seeds are $0,1,2$. Use 128 uniformly spaced training points and 1001 evaluation midpoints, which are distinct from those training points. Check endpoint initial data separately. Fix 5000 Adam steps with learning rate $10^{-3}$ followed by at most 200 L-BFGS iterations; report actual function evaluations and elapsed time, because identical iteration counts do not imply identical cost.

A high-accuracy DOP853 reference is checked against Radau; an exact formula is available for $\delta=0$. For cost comparison, vary DOP853 tolerances to obtain a practical error--cost curve rather than timing only the over-resolved reference. Report every seed, failures, derivative errors, $e_F$, $e_Y$, raw/normalized residuals, regularity margins and training/evaluation costs \cite{baselines}. A claimed speed benefit must specify accuracy, hardware, and whether training is amortized.

\section{Preparatory Verification Results}\label{sec:results}
The following computations were performed on 12 September 2026. They do not represent completion of the planned 2027 experiment.

\subsection{Symbolic and Classical Checks}
The verification package checks the arbitrary-profile identity~\eqref{eq:lift}, the exact power-law control, the printed traveling-wave balances, the Eulerian counterexample and a bounded source discrepancy: $\phi=ct+\Phi(s)$ satisfies the printed master equation for every sufficiently smooth positive-slope $\Phi$, so the quadratic restriction in source Equation (32) is unnecessary for that equation.

For the exact cubic control on 501 points, the maximum sampled state difference from the analytic formula is $\LagrangeControlError$. This near-roundoff observation is not a validated floating-point enclosure. For $\delta=0.01$, changing DOP853 step and tolerance settings gives difference $\LagrangeRefinement$; Radau differs by $\LagrangeOracle$. Sampled minima are $F'=\LagrangeMinR$ and $|a|=\LagrangeMinA$.

Five-point finite differences of the computed state curve, at step sizes $10^{-3}$ and $5\times10^{-4}$, check derivative consistency independently of the ODE right-hand side. This is a state-ODE diagnostic, not direct third differentiation of the interpolated position curve. The PDE lifting identity is separately symbolic. Sampled $P_2$ and $P_4$ profile errors for the perturbed case are $\LagrangeTaylorTwo$ and $\LagrangeTaylorFour$.

\subsection{Bounded Neural Feasibility Test}
To check operational feasibility, each neural trial was run with seed 0 for 400 Adam steps, without L-BFGS, on the perturbed profile. Table~\ref{tab:smoke} reports midpoint-grid maxima against the classical reference; the different evaluation grid explains the small difference from the preceding polynomial errors.
\begin{table}[htbp]
\centering
\caption{Preparatory single-seed CPU runs; no formal performance conclusion. Time covers training only.}\label{tab:smoke}
\begin{tabular}{ccccc}
\toprule
$N$ & Initial $e_F$ & Final $e_F$ & Final normalized defect & Seconds\\
\midrule
\tableinput proposal_support/feasibility_20260912/table.tex
\bottomrule
\end{tabular}
\end{table}

Both runs completed with finite gradients, zero computed initial-jet error and positive sampled slopes. Direct automatic differentiation of the original $(t,s)$ master expression on an $11\times11$ reconstruction grid agreed with the independently evaluated reduced-residual lift to below $3\times10^{-14}$. This verifies consistency of the implementations, not zero residual: the sampled master residuals remained approximately $0.477$ and $0.148$ for $N=2$ and $N=4$.

The $N=2$ run has lower profile error, whereas $N=4$ has lower normalized residual. The single-seed data therefore do not establish a higher-order advantage. The observed defect levels are also too large for the coarse full-interval factor in~\eqref{eq:certificate} to provide a useful containment certificate; uniform defect validation has not been performed. The exact box bound remains valid, while a certified trained-network error remains unestablished.

\section{Eight-Week Study in 2027}\label{sec:plan}
\begin{enumerate}
\item \textbf{Weeks 1--2: foundations and source audit.} Explain the mass coordinate, reproduce the selected symmetry and reduction, verify initial data and the retained model caveat.
\item \textbf{Weeks 3--4: classical approximation and theory.} Reproduce both IVPs, Taylor profiles and the conditional bound; tune practical classical baselines and document failed certificate conditions.
\item \textbf{Weeks 5--6: registered neural comparison.} Run both parameterizations on both profile cases for all three seeds with the declared schedule. Preserve raw metrics and failures.
\item \textbf{Weeks 7--8: analysis and write-up.} Produce error--cost comparisons, verify citations and reproducibility, and finalize the manuscript. Uniform validation of a trained network is a stretch task, not a graduation requirement.
\end{enumerate}
The core workload requires calculus, elementary ODE theory, one explicit Lie invariant and basic automatic differentiation. It does not require learning full differential elimination, training large operators or reconstructing the source's complete optimal system. Exact calendar dates within 2027 remain to be set.

\section{Limitations and Conclusion}\label{sec:conclusion}
The retained conclusion is specific: the printed master equation admits a directly verified cubic reduction with a computable regularity region, exact residual lifting and a conditional reduced-IVP certificate. A compact Taylor--PINN experiment is runnable on the available CPU. These facts justify a bounded undergraduate study; they do not establish a new general method or a neural advantage.

The physical conversion issue, restricted invariant family, singular leading coefficient, conditional uniform-bound premises, nonisolated Taylor/envelope comparison and incomplete multi-seed evidence remain explicit. A classical-solver advantage or an inconclusive certificate is an admissible research outcome. Established derivations and negative findings are preserved in the project claim register; stronger claims require additional checked evidence.

\appendix
\section{Exact Rational Bound Procedure}\label{app:interval}
Let $D=18\sigma^2r^3-5$. Differentiating~\eqref{eq:rationalf} gives
\begin{align}
f_r&=\frac{2(-54r^8\sigma^2+162r^6v^2\sigma^4+60r^5
+375r^4v\sigma+450r^3v^2\sigma^2-25v^2)}{r^2D^2},\\
f_v&=-\frac{2(25r^4\sigma+18r^3v\sigma^2+10v)}{rD}.
\end{align}
The code symbolically checks these against the derivative of $-b/a$.
Partition the three axes into 12 closed intervals each, with endpoints
\[
\sigma_i=1+i/24,\quad r_j=17/20+j/30,\quad v_k=-17/20+k/24,
\qquad i,j,k=0,\ldots,12.
\]
These 1728 subboxes cover the entire declared box, including boundaries.
Since $r>0$ and $D\ge24217/4000>0$, all interval divisions are valid.
Addition, multiplication, integer powers and reciprocals use exact rational endpoint bounds. On each subbox, add the absolute upper enclosures for $f_r$ and $f_v$ and take the maximum. The resulting rational number is the one in Proposition~\ref{prop:box}. Unit tests exercise sign-changing powers, negative reciprocals, zero-denominator rejection and the exact partition result. No claim is made that this proves the perturbed trajectory or a neural trial remains inside the box.

\section{Source Attribution and Reproducibility}\label{app:sources}
The publisher-hosted primary PDF was retrieved and found byte-identical to the supplied local file. Its SHA-256 is recorded in the companion source audit. The adopted source claims are limited to printed Equation (7), the selected cubic generator and invariant, the power-law root relation and the traveling-wave linearization theme. These are checked locally. The complete classification, optimal system and every expanded source ODE are not independently certified.

The methodological source audit records actual PDF sections, identity checks and access limits. Olver is cited as background, not through an invented theorem number. Taylor-model and symmetry-loss results are not generalized beyond their stated methods. Refined numerical references are not misattributed as forbidden by the baseline literature.

The companion claim register and source audit are:
\begin{quote}\small
\path{research_record.md}\\
\path{proposal_support/submission_source_audit.md}
\end{quote}
The executable entry points, run from the repository root, are:
\begin{quote}\small\ttfamily
python -B draft/proposal\_support/verify\_lagrange.py\\
python -B draft/proposal\_support/verify\_bounds.py\\
python -B draft/proposal\_support/pinn\_feasibility.py --steps 400\\
python -B -m unittest discover -s draft/proposal\_support
\end{quote}
The stored classical, exact-bound and neural reports include source hashes and evidence boundaries. Neural timing is machine-dependent. The primary PDE library and production symmetry engine are unchanged. Historical manuscripts and Eulerian near-cubic results remain archived and are not silently reused as current validation.

\input{proposal_support/references_2027.tex}
\end{document}
